\include{new_preamble}

\begin{document}
	
		\maketitle
	
	\newpage
	
	\begin{abstract}
		\noindent In the Hamming setup, we have seen that any code $\calc$ over the alphabet $\Sigma$ with distance $d$ can correct upto $\lfloor\frac{d-1}{2}\rfloor$ errors. The main goal of list decoding is to get past this bound. We shall also see that list decoding is, to an extent, the bridge between Hamming and Shannon.
		
		%% writing left %%% 
	\end{abstract}
	
	\newpage
	
	\section*{Notations}
	
	\noindent We denote a finite set of symbols, called an \emph{alphabet}, by $\Sigma$.
	\vspace{0.1in}
	
	\noindent For $n\in\mathbb{N} := \{1,2,\ldots\}$, we define $[n]$ to be the set $\{1,2,\ldots,n\}$.
	\vspace{0.1in}
	
	\noindent We denote the $q$-entropy function by $H_q(\rho) := \rho \log_q(q-1) - plog_qp - (1-p)log_q(1-p)$.
	\vspace{0.1in}
	
	\noindent Let $q$ be a prime power. Then $\mathbb{F}_q$ denotes the finite field of order $q$.
	\vspace{0.1in}
	
	\noindent We denote the Hamming distance by $\Delta(\bfx,\bfy) := \#\{i\in[n] | x_i \neq y_i\}$ for $\bfx, \bfy\in\Sigma^n$, and the Hamming Ball by $B(\bfy,e):=\{x\in\Sigma^n | \Delta(\bfx,\bfy)\leq e\}$.
	
	
	\newpage
	
	\tableofcontents	
	
	\newpage
	
	\section{Introduction}
	First we define what is a list decodable codes.
	\begin{definition}{}{}
		Given $0 \leq \rho <1$ and $L\geq1$, a code $\calc$ of block length $n$ over the alphabet $\Sigma$ is said to be $(\rho, L)$\emph{-list decodable} if for all $\bfy\in\Sigma^n$,
		\[
		\#\{\bfc\in\calc | \Delta(\bfy,\bfc)\leq\rho n\} \leq L.
		\]
	\end{definition}	
	Here this $\rho$ parameter is called \emph{fraction of error} and $L$ is called the \emph{list size upper bound}. Observe that, if fraction of error is atmost $\rho$, then the required codeword would be in the list. Also if $\calc$ is $(\rho, L)$-list decodable, then the decoding algortithm outputs atmost $L$ many codewords as outputs.
	
	\subsubsection*{Choice of $L$}
	Note that the case of $L=1$ is precisely the scenario of unique decoding, so it would be better to cosider $L>1$. Also the case $L=|\Sigma|^n$ make the notion of list decoding trivial. So we need to be careful about the choice of $L$.
	
	One way to restrict $L$ is to set to be an arbitrarily large constant. Now we can ask the following question.
	
	\begin{question}{}{}
		What is the limit of $\rho$ as $L\to\infty$ such that all sufficiently long codes $\calc\in\mathfrak{C}$ are $(\rho, L)$-list decodable ?
	\end{question}
	We refer to this as the list-decodability of $\mathfrak{C}$ with constant-sized lists. We can also see at a weker limit keeping our algorithmic goal in consideration. For an efficient algorithm, $L$ has to be a polynomial in $n$, as otherwise the output size would be super polynomial, giving us an algorithm not having a polynomial running time. This is the list decodability of $\mathfrak{C}$ with polynomial sized lists. For most of the natural codes, the above two notions are the same. For the rest of our discussion we shall only be interested in list decodability with polynomial sized lists.
	
	\subsubsection*{Utility of List Decoding}
	
	\begin{question}{}{first}
%		\label{first}
		Given a code of relative distance $\delta$, can we correct more than $\frac{\delta}{2}$ fraction of errors using list decoding?		
	\end{question}
	
	
	%% writing left %%%%
	
	
	\subsection{Johnson Bound}
	\begin{theorem}{Johnson Bound.}{johnson}
%		\label{johnson}
		If $\calc\subseteq[q]^n$ is a code of distance $d$ and $\rho<J_q(\frac{d}{n})$, then $\calc$ is $(\rho, qnd)$-list decodable, where the function $J_q(\delta)$ is defined as 
		\[
		J_q(\delta) = \left( 1 - \frac{1}{q}\right)\left(1 - \sqrt{1-\frac{q\delta}{q-1}}\right).
		\]
	\end{theorem}
	\emph{(The following proof is for the case of binary alphabet.)}
	
	\begin{proof}
		Denote the error $e = \rho n$. We have to study the code $\calc\subseteq\{0,1\}^n$ with distance $d$ and show that for all $\mathbf{y}\in\{0,1\}^n$, 
		\[
		|B(\mathbf{y},e)\cap \calc| \leq 2dn.
		\]
		Fix a code $\calc$ and $\mathbf{y}\in\{0,1\}^n$ and suppose $B(\mathbf{y}, e)\cap\calc = \{\mathbf{c_1},\ldots,\mathbf{c_M}\}$. So we need to show $M\leq 2nd$. Define $\mathbf{c_i'} = \mathbf{c_i}-\mathbf{y}$ for all $i\in[M]$. Consider the sum 
		\[
		S = \sum_{i<j}\Delta(\mathbf{c_i'},\mathbf{c_j'}).
		\]
		As the code is of distance $d$, then $\Delta(\mathbf{c_i'},\mathbf{c_j'})\geq d$ for all $i\neq j$, thus we get
		\begin{equation}
			\label{lb}
			S \geq \binom{M}{2}d.
		\end{equation}
		Consider the $n\times M$ matrix $\mathbf{A} = [\mathbf{c_1} : \cdots : \mathbf{c_M}]$. Define $m_i$ to be the number of $1$'s in the $i$-th row of $\mathbf{A}$. Observe that the $i$-th row contributes $m_i(M-m_i)$ to the sum $S$. Thus $S = \sum_im_i(M-m_i)$.  Define $\overline{e} := \frac{1}{M}\sum m_i$. Note that 
		\[
		\sum_{i=1}^n m_i = \sum_{j=1}^M wt(\mathbf{c_j'}) \leq eM
		\]
		as $wt(\mathbf{c_j'}) = \Delta(\mathbf{c_j}, \mathbf{y}) \leq e$. Thus $\overline{e}\leq e$. Also by \emph{RMS-AM} inequality, 
		
		\begin{align*}
			S = \sum_{i=1}^n m_i(M-m_i) 
			& = \overline{e}M^2 - \sum_{1}^{n}m_i^2 \\
			& \leq \overline{e}M^2 - \frac{(\overline{e}M)^2}{n} \\
			&  = M^2(\overline{e} - \frac{\overline{e}^2}{n}).
		\end{align*}
		
		So we get our combined inequality from the above inequality and (\ref{lb}), 
		\[
		\binom{M}{2} \leq M^2(\overline{e} - \frac{\overline{e}^2}{n}).
		\]
		
		This gives us an upper bound on $M$, which is the following, 
		\begin{align*}
			M \leq \frac{dn}{dn - 2n\overline{e} + 2\overline{e}^2} = \frac{2dn}{2dn - n^2 + n^2 - 4n\overline{e} + 4\overline{e}^2} 
			& = \frac{2dn}{(n - 2\overline{e})^2 - n(n-2d)}\\
			& \leq \frac{2dn}{(n-2e)^2 - n(n-2d)} \quad \text{as $\overline{e} \leq e$}. 
		\end{align*}	
		We also have the condition $\frac{e}{n}<J_2(\frac{d}{n})$, thus 
		\[
		\frac{e}{n} < \frac{1}{2} \left( 1- \sqrt{1-\frac{2d}{n}} \right)
		\Rightarrow n-2e > \sqrt{n(n-2d)}
		\Rightarrow (n-2e)^2 - n(n-2d) > 0.
		\]
		
		As all the terms there are natural numbers then, $(n-2e)^2 - n(n-2d) \geq 1$. Hence $M\leq 2dn$.
		
	\end{proof}
	
	This alone does not answer the Question~\ref{qsn:first}. Theorem~\ref{th:johnson} along with Lemma~\ref{lem:J_function_bound} gives us an affirmative answer to the Question~\ref{qsn:first}.
	
	\begin{lemma}{}{J_function_bound}
%		\label{J_function_bound}
		Let $q\geq 2$ be an integer and $0\leq x\leq 1-\frac{1}{q}$, then the following inequality holds.
		\[
		J_q(x) \geq 1 - \sqrt{1-x} \geq \frac{x}{2}.
		\]
		The second inequality is strict when $x\neq0$.
	\end{lemma}
	
	\begin{proof}
		Cosider the funciton $f(x) = J_q(x) - 1 + \sqrt{1-x}$. The derivative of this function is 
		\[
		f'(x) = \frac{1}{2}\left(\frac{1}{\sqrt{1 - \frac{qx}{q-1}}} - \frac{1}{\sqrt{1-x}}\right) \geq 0 \quad \text{for all $x$.}
		\]
		So the function is increasing and $f(0) = 0$, so we get our forst inequality. The second inequality is immediate from the inequality $(1-\frac{x}{2})^2\geq 1-x$.
	\end{proof}
	
	Theorem~\ref{th:johnson} and Lemma~\ref{lem:J_function_bound} imply that for a code of relative distance $\delta>0$, we correct strictly larger fraction of errors than unique decoding in ploynomial time as long as $q$ is at most some polynomial in $n$. These two propositions together also implies an \emph{alphabet free} version of the Johnson Bound.
	
	\begin{theorem}{Alphabet-Free Johnson Bound.}{}
		For any $q$-ary code $\calc$ of block length $n$ and distance $d$, if $e\leq n - \sqrt{n(n - d)}$, then $\calc$ is $(e/n, qnd)$-list decodable.
	\end{theorem}
	
	
	Now a natural question is 
	\begin{question}{}{}
		Is Johnson bound tight? Is it tight for linear codes?
	\end{question}
	
	We shall see there are codes, even linear codes, with relative distance $\delta$ so that we can find Hamming ball of radius greater than $J_q(\delta)$ with super polynomial many codewords. But we shall also see that, in some sense, it is not tight.
	
	\subsection{List Decoding Capacity}
	
	\begin{theorem}{}{ld_capacity}
%		\label{ld_capacity}
		For $q\geq2, 0\leq \rho < 1 - \frac{1}{q}$ and $\varepsilon>0$ small enough, there exist codes of block length $n$ large enough so that 
		\begin{enumerate}[1.]
			\item If $R\leq 1 - H_q(\rho) - \varepsilon$, then there is a $(\rho, O(\frac{1}{\varepsilon}))$-list decodable code.
			\item If $R\geq 1 - H_q(\rho) + \varepsilon$, then every $(\rho, L)$-list decodable code has $L\geq q^{\Omega(\varepsilon n)}$.
		\end{enumerate}  
	\end{theorem}
	
	Hence \emph{capacity of list decoding} is $1 - H_q(\rho)$. This is the exact capacity of a $qSC_{\rho}$ channel. Hence list decoding is the bridge between Hamming and Shannon.
	
	\subsubsection{Proof of Theorem~\ref{th:ld_capacity}}
	Rather than proving Theorem~\ref{th:ld_capacity} directly we prove the following  stronger theorem. In fact, we shall show that a random code is $(\rho, L)$-list decodable with high probability as long as 
	\[
	1 - H_q(\rho) - \frac{1}{L}.
	\]
	We shall show that probability of \emph{bad event} is small. \emph{Bad Event} means there exist $\mathbf{m_0},\ldots,\mathbf{m_L}\in [q]^{Rn}$ and $\mathbf{y}\in[q]^n$ so that 
	\[
	\Delta(\mathbf{C(m_i)}, \mathbf{y}) \leq \rho n \quad \text{for all $i\in \{0,\ldots,L\}$}.
	\]
	
	\begin{theorem}{List-Decoding Capacity.}{}
		Let $q\geq2$ be an integer, $0<\rho<1-\frac{1}{q}$ and $\varepsilon>0$.
		\begin{enumerate}[1.]
			\item Let $L\in\mathbb{N}$, then there exists a $(\rho, L)$-list decodable code with rate 
			\[
			R \leq 1 - H_q(\rho) - \frac{1}{L}.
			\]
			\item For all $(\rho, L)$-list decodable code with rate $R\geq 1 - H_q(\rho) + \varepsilon$, it is necesary that $L\geq q^{\Omega(\varepsilon n)}$.
		\end{enumerate}
	\end{theorem}
	\begin{proof}
		The following proof is done by Probabilistic Methods. 
		
		\vspace{0.01in}
		
\noindent \textsf{(Proof of $\mathsf{1}$).} Given a received word $\mathbf{y}\in[q]^n$, and $\mathbf{m_0}, \ldots, \mathbf{m_L}\in[q]^k$, the tuple $(\mathbf{y}, \mathbf{m_0}, \ldots, \mathbf{m_L})$ defines a \emph{bad event} if 
		\[
		C(\mathbf{m_i})\in B(\mathbf{y}, \rho n) \quad \text{for all $i\in\{0,\ldots,L\}$}.
		\]
		Note that $\calc$ is $(\rho, L)$-list decodable if there is no \emph{bad event}. 
		
		Fix $\mathbf{y}\in[q]^n$, and $\mathbf{m_0}, \ldots, \mathbf{m_L}\in[q]^k$. Note that for fixed $i$, by the choice of $\calc$, we have
		\[
		\Pr[C(\mathbf{m_i}) \in B(\mathbf{y}, \rho n)] = \frac{Vol_q(\rho n, n)}{q^n} \leq q^{-n(1 - H_q(\rho))}. 
		\]
		The probability of the \emph{bad event} by $(\mathbf{y}, \mathbf{m_0}, \ldots, \mathbf{m_L})$ is 
		\begin{align*}
			\Pr\left[\bigwedge_{i = 0}^L C(\mathbf{m_i})\in B(\mathbf{y}, \rho n)\right]  
			& =  \prod_{i = 0}^L \Pr[C(\mathbf{m_i})\in B(\mathbf{y}, \rho n)] \quad \text{[Indepence of $C(\mathbf{m_i})$'s]} \\
			& \leq q^{-n(L+1)(1 - H_q(\rho))}.
		\end{align*}
		
		Now we need to have no \emph{bad events} for $\calc$ in order to have it $(\rho, L)$-list decodable.
		\begin{align*}
			\Pr[\text{There is no bad event}]
			& \leq q^n\binom{q^k}{L+1}q^{-n(L+1)(1 - H_q(\rho))} \\
			& \leq q^n q^{k(L+1)} q^{-n(L+1)(1 - H_q(\rho))} \quad \text{[$\because \binom{n}{k} \leq n^k$]}\\ 
			& = q^{-n(L+1)[1 - H_q(\rho) - \frac{1}{L+1} - R]} \quad [R = kn] \\
			& \leq q^{-n(L+1)[1 - H_q(\rho) -\frac{1}{L+1} -1 + H_q(\rho) + \frac{1}{L}]} \\
			& = q^{-\frac{n}{L}} \\
			& < 1.
		\end{align*}
		The first inequality is apparent from going through all the received words and messages possible and using the union bound. Thus, by probabilistic method, there exists a code $\calc$ which is $(\rho, L)$-list decodable. Also the arguement tells us that probability of a random code not being $(\rho, L)$-list decodable goes to 0 as $n\to\infty$. \newline
		
	\noindent \textsf{(Proof of $\mathsf{2}$).} 
	\end{proof}
	
	
	\newpage
	
	\begin{thebibliography}{}
		
		\bibitem{guruswami}
		V.~Guruswami, A.~Rudra, and M.~Sudan, \textsf{Essential Coding Theory}, 
		Graduate Studies in Mathematics, vol.~226, American Mathematical Society, 2024.
	\end{thebibliography}
	
	
	
\end{document}
